% page 307 du fac-simile (image p-306 du scan)
% bloc 180.3 x 246.3 mm ; marge gauche 21.7 mm
\pgc{0.000mm}{0.000mm}{
\l{\cel{61}299}
\l{}
\l{\cel{5}\sou{konsiderar}. (\sou{trans)}.\cel{2}Atencar sorgoze, minucioze (ta o ca}
\l{\cel{8}motivo decidigiva, rezolvigiva) por agar saje. - DEFIRS.}
\l{}
\l{\cel{5}\sou{konsignar}. (\sou{trans)}. Donar (kozo konservenda kom garantiivo,}
\l{\cel{8}kom depozajo) : (l) pekunio-quanto, ye la kazo di tribunalo ;}
\l{\cel{8}(2) vari, depoze che komisionisto, komercisto - bagaji}
\l{\cel{8}depoze en relvoyo-staciono. - DEFIRS.}
\l{}
\l{\cel{5}\sou{konsilar}. (\sou{trans.,}\cel{1}ad). Guidar (ulu) per dicar a lu to, quon}
\l{\cel{8}lu devas agar o facar se lu volas agar saje. - EFIRS.}
\l{}
\l{\cel{5}\sou{konsistar}. (\sou{netrans., ye,}\cel{1}ek). Esar formacita (ek ula elementi).}
\l{\cel{8}- EFIS.}
\l{}
\l{\cel{5}\sou{konsistorio}. I. (\sou{che la katoliki)}. Asemblitaro de kardinali,}
\l{\cel{8}prezidata da la papo, qua su okupas per la aferi generala}
\l{\cel{8}di la eklezio katolika. - II. (\sou{che la protestantisti e}\cel{1}la}
\l{\cel{8}\sou{Izraelidi}). Asemblitaro ek kulto-oficieri e laiki, elektita}
\l{\cel{8}por direktar la aferi di la kongregaciono. - DEFIRS.}
\l{}
\l{\cel{5}\sou{konskriptar}. (\sou{trans.})\cel{2}En-armeigar koakte la adolecanti qui}
\l{\cel{8}jus atingis la evo-punto quan la lego determinis pri la}
\l{\cel{8}milito-kapabligo. - DEFIRS.}
\l{}
\l{\cel{5}\sou{konsolo}. I. (\sou{arkitekt.}) Saliajo fixigita an murego, por su-}
\l{\cel{8}portar balkono, kornico, e c., o kom piedestalo por statuo.}
\l{\cel{8}II. Orno-moblo pozita an muro, destinita a subtenar vazi,}
\l{\cel{8}kupi, statueti, e c. - DEF.}
\l{}
\l{\cel{5}\sou{konsolacar}. (\sou{trans., pri)}. Alejar la chagreno di (ulu). - EFIS.}
\l{}
\l{\cel{5}\sou{konsoldo}. (\sou{bot.)}\cel{2}Planto qua formacas genero di la familio}
\l{\cel{8}\textquotedbl{}boraginei\textquotedbl{}, di qua la radiko uzesis olim kom astriktivo. -}
\l{\cel{8}L.\cel{1}\sou{symphytum}.}
\l{}
\l{\cel{5}\sou{konsolidar}. (\sou{trans., per)}. (\sou{financo}).(\sou{aludante la debajo publik}a).}
\l{\cel{8}Transformar la debajo rimborsebla di stato, aden debajo perpe-}
\l{\cel{8}tua di qua nur la rento esos postulebla. - DEFIRS}
\l{}
\l{\cel{5}konsomeo. Buliono qua, pro koqueso longa-tempa, absorbis la}
\l{\cel{8}tota suko di la karno. - DEFIS.}
\l{}
\l{\cel{5}konsonancar. (netrans.) I. (aludante foni muzikala). Inter-}
\l{\cel{8}akordar plu o min komplete. - II. (\sou{aludante silabi, e pre}}
\l{\cel{8}cipue finali di qui la termino-foni esas sama). -Inter-}
\l{\cel{8}proximesar. - DEFIRS.}
\l{}
\l{\cel{5}konsonanto. Artikulo-maniero di la fono, qua varias segun la}
\l{\cel{8}movi da la lango o da la labii. (ex.: b, t, f esas konso-}
\l{\cel{8}nanti, kontre ke a, e, i\cel{2}esas vokali). - DEFIRS.}
\l{}
\l{\cel{5}konsorto.I. La persono di qua la estado sociala esas komuna kun}
\l{\cel{9}olta di sua spozo. (Dicesas precipue pri la spozulo di la}
\l{\cel{9}rejino di Britania.) - II.Singla de la personi qui, en afero,}
\l{\cel{8}havas interesto komuna. - DEFIRS.}
}
